% Copyright 2007 by Till Tantau
% Copyright 2014 by Rich Wareham, Filippo Spiga
% Copyright 2016 by Petar Veličković
% Copyright 2020 by Shivam N Patel
% This file may be distributed and/or modified:
% 1. under the LaTeX Project Public License and/or
% 2. under the GNU Public License.
%
% See the file LICENSE for more details.
% Adapted Spring 2026 by Sam Frederick

% ------------------------------------------------------------------------------
% STYLE NOTE:
% This custom template adds 'PartMC-themed' primary, secondary, and ternary 
% colors in ./cambridgecolors.sty and uses a warm off-white background color 
% 'WarmCream' defined in in beamercolorthemecambridge.sty. These can be adjusted
% to your liking.
% ------------------------------------------------------------------------------

\documentclass[aspectratio=169]{beamer}

\usetheme{cambridge}
\usefonttheme[onlymath]{serif}

\setbeamertemplate{navigation symbols}{}  % suppress navigation symbols
\setbeamercovered{transparent}

%\hypersetup{pdfpagemode=FullScreen,pdfpagelayout=SinglePage}

%%% Standard packages:
%\usepackage[english]{babel}
%\usepackage[latin1]{inputenc}
\usepackage{graphicx}
\usepackage{multicol}
\usepackage{subfigure}
\usepackage{listings}
%\usepackage[UKenglish]{isodate}

% Setup TikZ (if required)
%\usepackage{tikz}
%\usetikzlibrary{arrows}
%\tikzstyle{block}=[draw opacity=0.7,line width=1.4cm]

%\cleanlookdateon

% Custom packages
\usepackage{caption}
\usepackage[dvipsnames]{xcolor}
\usepackage[absolute,overlay]{textpos}

% Citations (if required)
\usepackage[style=apa,backend=biber]{biblatex}
\renewcommand*{\bibfont}{\small} % set small fontsize for reference slides
\setbeamertemplate{bibliography item}{} % remove document icons
\addbibresource{citations.bib}

% Author, Title, etc.
\title[Title abbreviation]
{%
 {\bf Your title here}%}
}

\author[Lastname, et al.]
{
  Firstname Lastname\inst{1}, Firstname Lastname\inst{2}
}

\institute[CL]
{\inst{1}Department of Climate, Meteorology, and Atmospheric Sciences, University of Illinois Urbana-Champaign,\\
\inst{2}Department of Mechanical Science and Engineering, University of Illinois Urbana-Champaign}

% Location and date of talk
\date[\today]
{Location of talk\hfill \today}

\begin{document}

\begin{frame}
  \titlepage
\end{frame}

% Sections appear in PDF reader navigation
%\section{The big idea}

\begin{refsection} % reference scope for overview

\begin{frame}{Aerosol thermodynamics: problem description}
  Gibbs free energy minimization is a \textbf{non-linear constrained optimization problem}. Specifically, the \textit{objective function} which we minimize (Gibbs free energy) is non-linear and non-convex with composition due to the non-linear activity coefficients \textcolor{Maroon}{$\gamma_i$}. Following \textcite{wexler_atmospheric_2002},
  \begin{align}
    \mathrm{minimize} \quad & G=\sum_i n_i\mu_i\\
    \mathrm{subject~to} \quad & \sum_i \nu_i n_i = b,\\
    \quad &  n_i \geq 0
  \end{align}
  Fortunately, the constraints (mass conservation and non-negativity), are linear. The chemical potential of each species $i$ is 
  \begin{equation}
    \mu_i = \mu_0(T, p) + RT\ln{\textcolor{Maroon}{\gamma_i} x_i}, 
  \end{equation}
  where the activity coefficient $\textcolor{Maroon}{\gamma_i}$ is a non-linear function of the multicomponent chemical composition. 
\end{frame}

% If you need a reference slide, otherwise comment out
\begin{frame}[allowframebreaks]{References}
  \printbibliography  
\end{frame} 
\end{refsection} % reference scope for overview

\end{document}